% ============================================================
%  FCFS Simulation: Main Conclusions by Mechanism
%  Meeting-ready 2-page conclusion sheet
%  May 2026
% ============================================================

\documentclass[11pt]{article}

\usepackage[margin=0.88in, top=0.78in, bottom=0.82in]{geometry}
\usepackage[T1]{fontenc}
\usepackage{amsmath, amssymb}
\usepackage{array}
\usepackage{booktabs}
\usepackage{enumitem}
\usepackage{mdframed}
\usepackage{microtype}
\usepackage{parskip}
\usepackage{tabularx}
\usepackage{titlesec}
\usepackage[table]{xcolor}
\usepackage[hidelinks]{hyperref}

\pagestyle{empty}
\setlength{\parindent}{0pt}
\setlength{\parskip}{0.35em}
\renewcommand{\arraystretch}{1.32}
\newcolumntype{Y}{>{\raggedright\arraybackslash}X}
\newcolumntype{L}[1]{>{\raggedright\arraybackslash\bfseries}p{#1}}

\titleformat{\section}{\large\bfseries}{}{0em}{}
\titlespacing*{\section}{0pt}{15pt}{6pt}

\setlist[itemize]{
  leftmargin=1.25em,
  itemsep=5pt,
  topsep=3pt,
  parsep=0pt,
  label=\textbullet
}

\newmdenv[
  backgroundcolor=gray!7,
  linecolor=gray!42,
  linewidth=0.6pt,
  roundcorner=2pt,
  innertopmargin=9pt,
  innerbottommargin=9pt,
  innerleftmargin=10pt,
  innerrightmargin=10pt,
  skipabove=8pt,
  skipbelow=10pt
]{infobox}

\newmdenv[
  backgroundcolor=blue!5,
  linecolor=blue!35!black,
  linewidth=0.7pt,
  roundcorner=2pt,
  innertopmargin=10pt,
  innerbottommargin=10pt,
  innerleftmargin=11pt,
  innerrightmargin=11pt,
  skipabove=8pt,
  skipbelow=10pt
]{recommendationbox}

\begin{document}

\begin{center}
  {\LARGE\bfseries FCFS Simulation: Main Conclusions by Mechanism}\\[5pt]
  {\normalsize May 2026 \(\cdot\) Two-class discrete-event FCFS appointment simulation}
\end{center}

\vspace{0.4em}

\begin{infobox}
This sheet summarizes the model-level conclusions from the FCFS simulation work.

The goal is to separate metric interpretation, behavioral mechanisms, and class-level effects.
\end{infobox}

\section*{1. Metric Conclusions}

\begin{itemize}
\item Utilization is not an access metric under oversubscription; it measures completed visits per slot.
\item Served rate is the primary aggregate access metric; it measures completed visits per arrival.
\item Offered wait must be reported with no-offer share because no-offer patients are absent from wait metrics.
\end{itemize}

\section*{2. Capacity Pressure}

\begin{itemize}
\item Increasing arrivals lowers served rate and raises offered wait when capacity is fixed.
\item At the stress-test baseline \((\lambda = 100,\ S = 32)\), high utilization \((0.841)\) coexists with poor access \((0.269)\).
\item The stress-test baseline is useful for mechanism diagnosis, not as a calibrated clinic estimate.
\end{itemize}

\section*{3. Behavioral Mechanisms}

\rowcolors{2}{gray!4}{white}
\begin{tabularx}{\textwidth}{L{0.22\textwidth} Y Y}
\toprule
Mechanism & Main effect & Interpretation \\
\midrule
Cancellation & Reduces own-class served rate. & May free future slots for other patients. \\
Balking & Reduces own-class served rate. & Reallocates access through the shared FCFS pool. \\
No-show & Directly lowers utilization. & Missed slots are not rebooked. \\
\bottomrule
\end{tabularx}

\newpage

\section*{4. Class-Level Conclusions}

\begin{itemize}
\item With symmetric parameters, the expected class access gap is near zero under pooled FCFS.
\item Class gaps emerge from asymmetric behavioral parameters or arrival rates, not explicit class priority.
\item In the realistic scenario, aggregate outcomes improve but the served-rate gap remains substantial \((0.220)\), so class-level metrics must be reported.
\end{itemize}

\section*{5. Reporting Recommendation}

\begin{recommendationbox}
{\large\bfseries Do not summarize the model with utilization alone.}

\vspace{0.5em}

Minimum reporting set:

\begin{itemize}
\item served rate;
\item utilization;
\item offered wait;
\item no-offer share;
\item outcome decomposition;
\item class-level served rates.
\end{itemize}
\end{recommendationbox}

\section*{6. Meeting Decision}

Meeting decision: choose the reference case for future analysis: stress-test baseline for mechanism diagnosis, or calibrated realistic scenario for operational interpretation.

\end{document}
